% ---- Preamble ----
% Check if the IsDraft command is defined
\ifdefined\IsDraft
  \documentclass[12pt, draft, twoside]{report} % maybe choose book?
\else
  \documentclass[12pt, twoside]{report} % maybe choose book?
\fi

\usepackage[utf8]{inputenc}
\usepackage{graphicx}
\graphicspath{ {images/} }
\usepackage[newlinetospace]{titlesec} % For title formatting
\usepackage{lipsum}
\usepackage[a4paper, width=125mm, top=45mm, bottom=45mm, bindingoffset=6mm]{geometry}
\usepackage{fancyhdr}
\pagestyle{fancy}
\fancyhead{}
% on the right for odd pages, on the left for even pages
\fancyhead[RO,LE]{\nouppercase{\leftmark}}
\renewcommand{\footrulewidth}{0pt}
\renewcommand{\headrulewidth}{1pt}
\setlength{\headheight}{15pt}

\ifdefined\IsRoboto
  \usepackage{roboto}
\fi

\usepackage[square]{natbib}
\bibliographystyle{plainnat}
\let\oldcite\cite
\renewcommand{\cite}[1]{\citep{#1}}

\usepackage{hyperref}
\hypersetup{
    colorlinks=true,
    linkcolor=black, % Making internal links black
    filecolor=black, % Making file links black
    urlcolor=blue, % Keeping URL links blue for distinction, but you can also set it to black
    citecolor=black, % Making citation links black
    anchorcolor=black, % Color for anchor text
    pdfborder={0 0 0}, % Removes the border around links
}


\usepackage[normalem]{ulem}
\usepackage{algorithm}
\usepackage{algpseudocode}
\usepackage{amsfonts}
\usepackage{amsmath}
\usepackage{amssymb}
\usepackage{amsthm}
\usepackage{array}
\usepackage{bbm}
\usepackage{bibentry}
\usepackage{booktabs}
\usepackage{breqn}
\usepackage{caption}
\usepackage{comment}
\usepackage{diagbox}
\usepackage{enumitem}
\usepackage{graphicx}
\usepackage{latexsym}
\usepackage{listings}
\usepackage{mathtools}
\usepackage{multirow}
\usepackage{physics}
\usepackage{subcaption}
\usepackage{textcomp}
\usepackage{thm-restate}
\usepackage{tikz-cd}
\usepackage{times}
\usepackage{url}
\usepackage{xcolor}
\usepackage{xspace}

\nobibliography*
\renewcommand{\norm}[1] { \left\Vert #1 \right\Vert }
\newcommand{\defeq}{\vcentcolon=}
\newcommand{\eqdef}{=\vcentcolon}
\newcommand{\prob}{\operatorname*{Pr}}
\DeclareMathOperator*{\argmax}{arg\,max}
\DeclareMathOperator*{\argmin}{arg\,min}
\newtheorem{theorem}{Theorem}
\newtheorem{corollary}{Corollary}
\newtheorem{proposition}{Proposition}
\newtheorem{definition}{Definition}
\newtheorem{lemma}[theorem]{Lemma}
\newcommand\la{\textlangle\xspace}
\newcommand\ra{\textrangle\xspace}

% ---- Document ----

\begin{document}

% ---- Title Page ----

\begin{titlepage}
{
\ifdefined\IsRoboto
\fontfamily{Roboto-TLF}\selectfont
\fi
    
   \centering
   \vspace*{-3.5cm}
   \includegraphics[width=0.5\textwidth]{sns_logo.pdf} 
   \par\vspace{2.5cm}

   {\Large Classe di Scienze\par}
   \vspace{0.5cm}
   {\Large Corso di perfezionamento in\par}
   {\Large \la nome-del-corso \ra\par}
   \vspace{0.5cm}
   {\Large \la numero-ciclo ciclo \ra \par}
   \vfill
   {\huge\bfseries \itshape \la Titolo: \ra \par}
   {\Large\bfseries \itshape \la Sotto-titolo \ra \par}
   \vspace{1.5cm}
   {\Large Settore Scientifico Disciplinare \la INF/01 \ra \par} % oppure ING-INF/05 ??
   \vfill
   \raggedright
   {\Large Candidato\par}
   {\Large \la nome \ra \par}
   \vspace{1.5cm}
   %\vfill
   {\large Relatore \par}
   {\large Prof. \la nome \ra \par}
   \vfill
   \centering
   {\large Anno accademico \la 2023/2024 \ra \par}
   \vspace*{-2.5cm}
   }
\end{titlepage}

% ---- Chapters ----

\pagenumbering{roman}

\chapter*{Abstract}
This is an abstract.

\listoffigures

\listoftables

\tableofcontents
\newpage

\pagenumbering{arabic}

\chapter{Introduction} \label{ch:intro}
This is an introduction.

\chapter{Conclusion} \label{ch:conclusions}
These are the conclusions.

\bibliography{bibliography}

\end{document}
